



The following command creates the Reed Muller code $\mbox{RM}_{3}.$ 
The command exports the code to magma and writes a file of all codewords.

\sourcecode{RM_3_export}

Here is the csv file containing the codewords. 

\orbiteroutput{GRAPHICS/WRAPPERS/RM3_codewords}

The four columns represent:
\begin{enumerate}
\item
Idx -- The numerical value of the message.
\item
Message -- The message.
\item
Codeword -- The associated codeword.
\item
CodewordRank -- The numerical value of the codeword.
\end{enumerate}

\bigskip



The next command creates a diagram of the codewords in the appropriate Hamming space, 
which is $H(8,2)$ in this case.

\sourcecode{RM_3_Hamming_space_diagram}

The following command produces a diagram of the characteristic 
function of the code in the Hamming space $H(8,2).$

\sourcecode{RM_3_draw}

The following boolean function representation of ${\rm RM}_{3}$ in $H(8,2)$ is produced:

\orbiteroutput{GRAPHICS/WRAPPERS/RM_3_diagram_01_8_16_draw}

The next command splits the vectors according to their distance from the code:

\sourcecode{RM_3_split}

We produce a diagram of each distance set. 
We loop over all distances $\delta=0,1,2$:

\sourcecode{RM_3_distance_draw}

The distance sets are shown below
\orbiteroutput{GRAPHICS/WRAPPERS/RM_3_distance_sets}
For each distance, 
we compute the algebraic normal form of the characteristic function of the distance set:

\sourcecode{RM_3_algebraic_normal_form}



\bigskip

In the example, we will compute a graphical representation of the codewords 
of the Reed-Muller code of order 6. First, we use makefile variables to 
define the rows of the generator matrix:

\sourcecodevariable{RM_6_GENERATOR_1}


\sourcecodevariable{RM_6_GENERATOR_2}


\sourcecodevariable{RM_6_GENERATOR_3}


\sourcecodevariable{RM_6_GENERATOR_4}


\sourcecodevariable{RM_6_GENERATOR_5}


\sourcecodevariable{RM_6_GENERATOR_6}


\sourcecodevariable{RM_6_GENERATOR_7}

The following command creates the generator matrix, and exports the generator matrix 
and the set of codewords to two csv files. 

\sourcecode{RM_6}

The next command produces a drawing of the generator matrix:

\sourcecode{RM_6_genma_draw}


The next command produces individual pictures for each of the 128 codewords. 
The command utilizes the \verb'-loop' command to automatize the processing. 
Afterwards, the codeword pictures are composed in one poster, 
using the \verb'magick' tool on the command line.

\sourcecode{RM_6_codewords_draw}

Finally, the 128 codewords are stitched together in the form of a tapestry.
Each codeword is a bitvector of length 64, rearranged as $8 \times 8$ bitmatrix.
\orbiteroutput{GRAPHICS/WRAPPERS/poster_RM_1_6_draw}


\bigskip


